\documentclass[11pt]{article}

\usepackage{amsmath,amssymb}
\usepackage{xurl}
\usepackage[colorlinks=true,linkcolor=blue,citecolor=blue,urlcolor=blue]{hyperref}

% Shared commands for notes in docs/paper-gaps/.

\DeclareUnicodeCharacter{2080}{\ifmmode{}_{0}\else\textsubscript{0}\fi}
\DeclareUnicodeCharacter{2081}{\ifmmode{}_{1}\else\textsubscript{1}\fi}
\DeclareUnicodeCharacter{2082}{\ifmmode{}_{2}\else\textsubscript{2}\fi}
\DeclareUnicodeCharacter{2099}{\ifmmode{}_{n}\else\textsubscript{n}\fi}

\newcommand{\Lean}{\textsc{Lean}}
\ExplSyntaxOn
\NewDocumentCommand{\leanid}{m}{
  \ifmmode
    \textsf{\detokenize{#1}}
  \else
    \group_begin:
    \urlstyle{sf}
    \tl_set:Nn \l_tmpa_tl {#1}
    \tl_replace_all:Nnn \l_tmpa_tl {\_} {_}
    \exp_args:NV \path \l_tmpa_tl
    \group_end:
  \fi
}
\ExplSyntaxOff
\providecommand{\tr}{\operatorname{tr}}
\newcommand{\tnleanIssueBase}{https://github.com/LionSR/TNLean/issues/}
\newcommand{\tnleanPrBase}{https://github.com/LionSR/TNLean/pull/}
\newcommand{\tnleanDocBase}{https://sirui-lu.com/TNLean/docs/}
\newcommand{\ghissue}[1]{\href{\tnleanIssueBase#1}{GitHub issue \##1}}
\newcommand{\ghpr}[1]{\href{\tnleanPrBase#1}{PR \##1}}
\newcommand{\doclink}[1]{\href{\tnleanDocBase#1.html}{\path{#1}}}

% Dirac notation and the identity operator, as in blueprint/src/macros/common.tex.
% \providecommand keeps note-local definitions authoritative.
\usepackage{dsfont}
\providecommand{\ket}[1]{|#1\rangle}
\providecommand{\bra}[1]{\langle #1|}
\providecommand{\braket}[2]{\langle #1 | #2 \rangle}
\providecommand{\ketbra}[2]{|#1\rangle\!\langle #2|}
\providecommand{\Id}{\mathds{1}}
\providecommand{\one}{\mathds{1}}
\providecommand{\C}{\mathbb{C}}
\providecommand{\MN}[1]{M_{#1}(\C)}
\providecommand{\MD}{M_D(\C)}

% External referencing for paper sources.  The build helper writes sanitized
% auxiliary files under build/paper-aux/ and excludes duplicated source labels.
\usepackage{xr}

% Import a generated paper auxiliary file with a caller-supplied label prefix.
% The candidate paths support builds from docs/paper-gaps/, the repository root,
% and build/paper-aux/ itself.
\newcommand{\importpaperaux}[2]{%
  \IfFileExists{../../build/paper-aux/#2.aux}{%
    \externaldocument[#1]{../../build/paper-aux/#2}%
  }{%
    \IfFileExists{build/paper-aux/#2.aux}{%
      \externaldocument[#1]{build/paper-aux/#2}%
    }{%
      \IfFileExists{#2.aux}{%
        \externaldocument[#1]{#2}%
      }{}%
    }%
  }%
}

% General external reference macros
\newcommand{\paperref}[2]{\ref{#1-#2}}
\newcommand{\papereqref}[2]{(\ref{#1-#2})}

% CPSV16 source (1606.00608 / MPDO-22-12-17-2.tex) external document and shortcuts
\importpaperaux{cpsv16-}{1606.00608/MPDO-22-12-17-2-tnlean}
\newcommand{\cpsvref}[1]{\paperref{cpsv16}{#1}}
\newcommand{\cpsveqref}[1]{\papereqref{cpsv16}{#1}}

% Verdict marker: \gapnote{<kind>}{<status>}. Kind is the policy
% classification (clarification, local-correction, scope-restriction,
% unfaithful, false-source, open-gap); status is open, wip (elimination
% underway), resolved, or historical. Machine-read by the site index;
% prints a verdict line.
\newcommand{\gapnote}[2]{\par\noindent\textbf{Verdict:} #1 (#2).\par}


\title{Complete Positivity in the Projected Support Corner Star-Algebra}
\author{}
\date{\today}

\begin{document}
\maketitle

\section{The assertion}

Corollary 6.6 of \emph{Quantum Channels \& Operations} (Wolf 2012), stated at
lines 1380--1510 of \texttt{Notes/WolfNoteTexSource/ch06\_spectral\_properties.tex},
concerns a positive trace-preserving linear map $T$ on $M_D(\mathbb{C})$ whose trace
adjoint $T^{*}$ is unital and satisfies the Schwarz inequality. For a positive
semidefinite fixed point $\rho$ of $T$ with support projection $Q$, the corner-restricted
fixed-point set
\[
    \{\, Y \in Q\,M_D(\mathbb{C})\,Q \;\mid\; Q\,T^{*}(Y)\,Q = Y \,\}
\]
is a $*$-subalgebra of the corner algebra $Q\,M_D(\mathbb{C})\,Q$.

Wolf offers complete positivity as an \emph{example} of a map meeting the hypothesis,
not as part of the hypothesis. The Schwarz inequality for the adjoint is strictly weaker:
every completely positive unital map satisfies it, but positive maps that are not
completely positive can satisfy it as well.

\section{The restriction in the formalization}

The available Lean statement is \leanid{Kraus.cornerFixedPointsStarSubalgebra}, with the
membership characterization \leanid{Kraus.mem\_cornerFixedPointsStarSubalgebra}, both in
\texttt{TNLean/Channel/FixedPoint/CornerFixedPoints.lean}. Its hypotheses are a family
$K$ of Kraus operators with \leanid{IsTP} $K$ and a positive semidefinite fixed point of
the induced map. Presenting the map by a Kraus family assumes complete positivity, so the
Lean hypothesis set is strictly stronger than the source's. Under the repository
faithfulness rule this declaration establishes the completely positive case of
Corollary 6.6 and is not a formalization of the corollary itself.

The corresponding Blueprint entry states the Kraus form explicitly in its hypotheses, so
its \leanid{\textbackslash leanok} is honest for the restricted statement; it must not be
read as the unrestricted corollary. The analogous restriction for Corollary 6.7 is
recorded separately in \texttt{wolf\_cor67\_maximal\_support\_restriction.tex}.

\section{Elimination plan}

No further material is needed from the source: the corollary is stated in full at the
lines cited above, and the machinery to express its hypothesis without Kraus data already
exists. The multiplicative domain of a bare linear map
$E : M_D(\mathbb{C}) \to M_D(\mathbb{C})$ is available as \leanid{multiplicativeDomain} in
\texttt{TNLean/Channel/Schwarz/AbstractMultiplicativeDomain.lean}, together with its
membership characterization, so the Schwarz hypothesis can be carried directly.

The elimination is therefore to restate the corollary for a positive trace-preserving map
whose trace adjoint is unital and Schwarz, discharging it through the general support
compression and the now-available Theorem 6.14, and then to recover the present Kraus
declarations as specializations rather than maintaining a parallel development. The
source-general statement is the one that may carry a source-facing Blueprint entry for
Corollary 6.6.

\end{document}
